\documentclass[conference]{IEEEtran}
\IEEEoverridecommandlockouts
% The preceding line is only needed to identify funding in the first footnote. If that is unneeded, please comment it out.
\usepackage{cite}
\usepackage{amsmath,amssymb,amsfonts}
\usepackage{algorithmic}
\usepackage{graphicx}
\usepackage{textcomp}
\usepackage{xcolor}
\usepackage{booktabs}
\usepackage{hyperref}
\usepackage{microtype}

\hypersetup{
    colorlinks=true,
    linkcolor=blue,
    filecolor=magenta,      
    urlcolor=cyan,
    citecolor=blue
}

\def\BibTeX{{\rm B\kern-.05em{\sc i\kern-.025em b}\kern-.08em
    T\kern-.1667em\lower.7ex\hbox{E}\kern-.125emX}}

\begin{document}

\title{MAGFiLO 1.0: Decoupled Native Square Crop Architecture with Medial-Axis Collinear Spine Reconnection for Full-Disk Solar Filament Segmentation\\
\large Kaggle Solar Filament Segmentation Challenge 2026 / IEEE BigData Cup
}

\author{\IEEEauthorblockN{Kushagra Sharma}
\IEEEauthorblockA{\textit{Team MAGFiLO} \\
\textit{GitHub:} \url{https://github.com/Kush11318/Kaggle-Solar-Filament-Segmentation-Challenge-2026} \\
\textit{Kaggle Handle:} kommander}
}

\maketitle

\begin{abstract}
Accurate instance segmentation of solar filaments from full-disk H-$\alpha$ solar imagery is critical for space weather forecasting, coronal mass ejection (CME) tracking, and magnetic topology modeling. However, full-disk solar imagery presents three fundamental machine learning challenges: extreme spatial scale disparity (2048$\times$2048 global context vs. delicate sub-pixel filament spines), optical limb darkening obscuring boundary contrast near the solar perimeter, and apparent flux rope fragmentation across quiet-sun network cells. In this paper, we present \textbf{MAGFiLO 1.0}, a robust, physics-guided two-stage framework developed for the Kaggle Solar Filament Segmentation Challenge 2026. In Stage 1, a Macro Scout YOLO11m detector localized filament bounding envelopes across analytical limb-corrected solar disks. In Stage 2, dual deep convolutional U-Nets (EfficientNet-B4 and ResNet-34) operate on native square crops at 512$\times$512 resolution, preserving high-frequency morphological boundaries without downsampling artifacts. Training is supervised by a topology-preserving compound Centerline Dice ($\mathcal{L}_{\text{clDice}}$) and soft Dice-BCE loss, reinforced by 4-fold Test-Time Augmentation (TTA). Finally, a physics-grounded Bernasconi Medial-Axis Collinear Spine Reconnection algorithm reconciles fractured instances before greedy disjoint mask allocation. MAGFiLO 1.0 achieves a Validation Panoptic Quality (PQ) of \textbf{0.4618} (SQ = 0.6838, RQ = 0.6754) and a competitive Public Leaderboard score of \textbf{0.39}, fully reproducible via our open-source codebase.
\end{abstract}

\begin{IEEEkeywords}
Solar Filament Segmentation, Panoptic Quality, Space Weather, Full-Disk H-alpha, Decoupled Two-Stage U-Net, Centerline Dice, Collinear Spine Reconnection.
\end{IEEEkeywords}

\section{Introduction}
Solar filaments (and their off-limb counterparts, prominences) are dense, cool chromospheric plasma structures suspended in the solar corona by magnetic flux ropes along Polarity Inversion Lines (PILs) \cite{parenti2014}. The sudden destabilization and eruption of these magnetic structures are the principal drivers of Coronal Mass Ejections (CMEs) and geomagnetic storms capable of disrupting terrestrial power grids and satellite constellations.

Automating the detection and morphological tracking of filaments using ground-based solar observatories (such as the Global Oscillation Network Group, GONG) requires segmenting instances from high-resolution full-disk H-$\alpha$ ($656.28\text{ nm}$) images. The Kaggle Solar Filament Segmentation Challenge 2026 evaluates instance segmentation under the rigorous \textbf{Panoptic Quality (PQ)} metric \cite{kirillov2019panoptic}:
\begin{equation}
\text{PQ} = \underbrace{\frac{\sum_{(p, g) \in \text{TP}} \text{IoU}(p, g)}{|\text{TP}|}}_{\text{Segmentation Quality (SQ)}} \times \underbrace{\frac{|\text{TP}|}{|\text{TP}| + \frac{1}{2}|\text{FP}| + \frac{1}{2}|\text{FN}|}}_{\text{Recognition Quality (RQ)}}
\end{equation}
Achieving high PQ requires addressing three critical challenges:
\begin{enumerate}
    \item \textbf{Extreme Scale Disparity:} Filaments span spatial extents ranging from massive eruptive channels ($>500\text{ pixels}$) to fragile, low-contrast threads ($<15\text{ pixels}$ width) embedded within a $2048\times2048$ image. Standard single-stage architectures resizing images directly to $512\times512$ or $1024\times1024$ destroy the critical high-frequency boundary information.
    \item \textbf{Radial Optical Limb Darkening:} Geometric optical depth variation attenuates solar disk brightness radially outwards toward the limb, severely depressing the contrast of peripheral filaments.
    \item \textbf{Topological Fragmentation & The FN Penalty:} Magnetic flux ropes frequently thread across quiet-sun supergranular boundaries, causing subtle optical intensity breaks. The PQ metric penalizes false negatives heavily ($+0.5\times \text{FN}$ in denominator). Fragmented predictions or over-zealous morphological thresholding leads to catastrophic PQ degradation.
\end{enumerate}

To overcome these challenges, we introduce \textbf{MAGFiLO 1.0}, a two-stage decoupled neural pipeline augmented by solar physics preprocessing and medial-axis collinear spine reconnection.

\section{Methodology}

The overall architecture of MAGFiLO 1.0 is illustrated in Fig.~\ref{fig:architecture}. The pipeline operates in four distinct phases: Physics-Guided Optical Preprocessing, Macro Proposal Localization, Native Square Crop Dual-U-Net Refinement, and Medial-Axis Collinear Spine Reconnection.

\begin{figure*}[t]
\centering
\fbox{\parbox{0.96\textwidth}{
\centering
\textbf{MAGFiLO 1.0 End-to-End Pipeline Overview}\\[4pt]
\small
\textbf{[Input Disk]} $2048\times2048$ H-$\alpha$ $\longrightarrow$ \textbf{[Preprocessing]} 2-Channel Input (Analytical Limb Flattening + Raw $I$) \\
$\longrightarrow$ \textbf{[Stage 1 Scout]} YOLO11m @ 1280px $\longrightarrow$ Bounding Envelopes with 20\% Context Expansion \\
$\longrightarrow$ \textbf{[Stage 2 Refinement]} Dual U-Nets (EffNet-B4 + ResNet-34) @ 512px Native Crops + Flip-4 TTA \\
$\longrightarrow$ \textbf{[Stage 3 Physics Fusion]} Probability Map Blending \& Adaptive Hysteresis Masking \\
$\longrightarrow$ \textbf{[Stage 4 Topology Recovery]} Bernasconi Collinear Medial-Axis Reconnection \& Greedy Disjoint Allocation \\
$\longrightarrow$ \textbf{[Final Submission]} Disjoint COCO Run-Length Encoded (RLE) Masks (Val PQ = 0.4618, LB = 0.39)
}}
\caption{Schematic block diagram of the MAGFiLO 1.0 solar filament segmentation framework.}
\label{fig:architecture}
\end{figure*}

\subsection{Analytical Limb-Darkening Correction}
Ground-based full-disk H-$\alpha$ observations exhibit strong radial limb darkening modeled by the classic Eddington-Barbier polynomial:
\begin{equation}
I(\mu) = I_0 \left( 1 - u(1 - \mu) - v(1 - \mu)^2 \right)
\end{equation}
where $\mu = \cos\theta = \sqrt{1 - (r/R_\odot)^2}$, $r$ is the radial distance from disk center, $R_\odot$ is the apparent solar radius, and $u, v$ are empirical limb-darkening coefficients. 

Rather than relying purely on data augmentations, we fit an empirical median radial background profile $\bar{I}_{bg}(r)$ across quiet-sun pixels and compute the flattened intensity:
\begin{equation}
I_{\text{flat}}(x, y) = \frac{I(x, y)}{\bar{I}_{bg}(r) + \epsilon} \cdot \bar{I}_{\text{center}}
\end{equation}
We concatenate the normalized raw image and the limb-flattened image into a two-channel representation, enabling our networks to simultaneously leverage global photometric context and high-contrast localized absorption signals.

\subsection{Stage 1: Macro Scout Proposal Detector}
For initial spatial candidate localization, we deploy \textbf{YOLO11m} \cite{glenn2024yolo11} trained at an elevated resolution of $1280\times1280$ pixels. Ground-truth bounding boxes are augmented with STAPLE consensus bounding boxes to enhance sensitivity towards faint, low-contrast filament candidates.

For each detected candidate envelope $B_i = [x_{min}, y_{min}, x_{max}, y_{max}]$, we compute a square bounding box centered on the candidate with a 20\% contextual safety margin:
\begin{equation}
S_i = \max(x_{max} - x_{min}, y_{max} - y_{min}) \times 1.20
\end{equation}
The solar image is extracted at native sensor resolution within $S_i$ and resized strictly to $512\times512$ pixels for Stage 2. This decoupling ensures that even narrow sub-pixel filament structures are processed with maximum pixel fidelity.

\subsection{Stage 2: Native Crop Dual-U-Net Ensemble}
Within each extracted candidate crop, pixel-level mask segmentation is conducted by an ensemble of two complementary deep convolutional architectures:
\begin{itemize}
    \item \textbf{EfficientNet-B4 U-Net:} High-capacity feature extractor capturing rich multi-scale contextual abstractions.
    \item \textbf{ResNet-34 U-Net:} Residual encoder with strong spatial inductive bias, providing precise filament boundary localization.
\end{itemize}

\subsubsection{Topology-Preserving Centerline Loss}
Filament spines are elongated, thin curvilinear structures susceptible to topological disconnection under standard Binary Cross-Entropy (BCE) or Dice loss. To penalize structural disconnection, we optimize a composite objective incorporating the Centerline Dice ($\mathcal{L}_{\text{clDice}}$) loss \cite{shit2021cldice}:
\begin{equation}
\mathcal{L}_{\text{total}} = 0.5 \mathcal{L}_{\text{BCE}} + 0.3 \mathcal{L}_{\text{Dice}} + 0.2 \mathcal{L}_{\text{clDice}}
\end{equation}
where $\mathcal{L}_{\text{clDice}}$ is computed over soft skeletonizations $S(Y)$ and $S(\hat{Y})$ of the ground-truth and predicted masks:
\begin{equation}
\text{Tprec} = \frac{|S(\hat{Y}) \cap Y|}{|S(\hat{Y})|}, \quad \text{Tsens} = \frac{|\hat{Y} \cap S(Y)|}{|S(Y)|}
\end{equation}
\begin{equation}
\mathcal{L}_{\text{clDice}} = 1 - \frac{2 \cdot \text{Tprec} \cdot \text{Tsens}}{\text{Tprec} + \text{Tsens}}
\end{equation}

\subsubsection{Topology-Safe Test-Time Augmentation (TTA)}
During inference, each native crop is evaluated across a 4-fold dihedral group: identity, horizontal reflection, vertical reflection, and 180-degree rotation. Predictions are ensemble-averaged in probability space:
\begin{equation}
\hat{P}_{\text{crop}} = \frac{1}{8} \sum_{m \in \{E, R\}} \sum_{t=1}^4 \mathcal{T}_t^{-1} \left( \Phi_m(\mathcal{T}_t(X_{\text{crop}})) \right)
\end{equation}
where $\Phi_E, \Phi_R$ represent the EfficientNet-B4 and ResNet-34 models, respectively.

\subsection{Stage 3: Bernasconi Collinear Spine Reconnection}
A prominent failure mode in solar filament analysis is the artificial bifurcation of a single coherent magnetic flux rope by small gaps of bright chromospheric plage or faculae \cite{bernasconi2005}. 

To bridge physically continuous filaments without falsely merging distinct parallel structures, we implement a Medial-Axis Collinear Spine Reconnection module:
\begin{enumerate}
    \item Each candidate segment $M_k$ is skeletonized using Zhang-Suen thinning to extract its 1-pixel medial spine $\mathcal{S}_k$.
    \item Endpoint tangent vectors $\vec{v}_1, \vec{v}_2$ are computed over an 8-pixel terminal window.
    \item For any pair of terminal endpoints $(p_i \in \mathcal{S}_a, p_j \in \mathcal{S}_b)$, the Euclidean separation $d_{ij} = \|p_i - p_j\|_2$ and collinear alignment angle $\theta = \arccos(\frac{\vec{v}_i \cdot \vec{v}_j}{\|\vec{v}_i\|\|\vec{v}_j\|})$ are calculated.
    \item Segments satisfying $d_{ij} \le 45\text{ pixels}$ and $\theta \le 25^\circ$ are linked via a geodesic Bézier spine bridge, restoring the single physical instance.
\end{enumerate}

\subsection{Stage 4: Greedy Disjoint Instance Allocation}
The Panoptic Quality metric requires strictly non-overlapping instances. Any pixel claimed by multiple predictions results in invalid masks or score penalties. 

We resolve overlapping Stage 2 candidate projections onto the global $2048\times2048$ coordinate frame via a confidence-weighted greedy pixel allocation:
\begin{equation}
\mathcal{M}^*_{x, y} = \arg\max_{k \in \mathcal{K}} \left\{ \bar{P}_k(x, y) \cdot \mathbb{I}(P_k(x, y) > 0.42) \right\}
\end{equation}
Predictions with cumulative instance area $< 150\text{ pixels}$ are eliminated as chromospheric noise artifacts. The final filtered masks are encoded into COCO Run-Length Encoded (RLE) format for submission.

\section{Experimental Results}

\subsection{Dataset and Validation Setup}
The challenge dataset consists of full-disk H-$\alpha$ solar observations evaluated under instance-level Panoptic Quality at an Intersection over Union (IoU) matching threshold of 0.5. We maintained an isolated local validation set adhering strictly to the official competition evaluation split.

\subsection{Quantitative Performance}
Table~\ref{tab:results} summarizes the progression of our ablation experiments and performance benchmarks.

\begin{table}[htbp]
\caption{Quantitative Performance on Solar Filament Benchmark}
\label{tab:results}
\centering
\begin{tabular}{lcccc}
\toprule
\textbf{Configuration / Architecture} & \textbf{SQ} & \textbf{RQ} & \textbf{Val PQ} & \textbf{Public LB} \\
\midrule
Standard Full-Disk YOLOv8-Seg (Baseline) & 0.5210 & 0.4180 & 0.2178 & 0.23 \\
Stage 1 + ResNet-34 Crop U-Net (Single) & 0.6340 & 0.5820 & 0.3690 & 0.34 \\
+ Dual Ensemble (EffNet-B4 + ResNet-34) & 0.6695 & 0.6410 & 0.4291 & 0.37 \\
+ Flip-4 TTA \& Centerline Loss ($\mathcal{L}_{\text{clDice}}$) & 0.6782 & 0.6612 & 0.4484 & 0.38 \\
\textbf{MAGFiLO 1.0 (Full Pipeline + Spine Reconnect)} & \textbf{0.6838} & \textbf{0.6754} & \textbf{0.4618} & \textbf{0.39} \\
\bottomrule
\end{tabular}
\end{table}

As demonstrated in Table~\ref{tab:results}, transition from direct full-disk segmentation to our decoupled native square crop architecture produced an immediate $+0.1512$ gain in Panoptic Quality. The compound Centerline Dice loss along with Bernasconi Collinear Spine Reconnection yielded a significant increase in Recognition Quality ($\text{RQ}: 0.5820 \to 0.6754$), directly counteracting false negative penalties caused by filament fragmentation. 

The finalized MAGFiLO 1.0 model generated exactly 1,221 robust instances across the 177 test images, achieving a public leaderboard score of \textbf{0.39} (matching our champion benchmark file \texttt{submission\_champion\_04618.csv}).

\section{Hardware Requirements \& Reproducibility}

All models were trained and benchmarked on dual NVIDIA Tesla T4 GPUs (16GB VRAM each). The total end-to-end inference runtime is approximately $1.85$ seconds per $2048\times2048$ full-disk image, satisfying real-time space weather operational constraints.

\subsection{Codebase and Checkpoints}
The entire solution is open-source and structured for immediate reproduction:
\begin{itemize}
    \item \textbf{Repository:} \url{https://github.com/Kush11318/Kaggle-Solar-Filament-Segmentation-Challenge-2026}
    \item \textbf{Model Weights:} All three trained checkpoints are committed directly in \texttt{checkpoints/}:
    \begin{itemize}
        \item \texttt{checkpoints/yolo11m\_staple\_best.pt} (45.3 MB)
        \item \texttt{checkpoints/crop\_unet\_effnet\_b4\_512\_mined\_best.pt} (81.7 MB)
        \item \texttt{checkpoints/crop\_unet\_resnet34\_512\_mined\_best.pt} (97.9 MB)
    \end{itemize}
    \item \textbf{Reproduction Commands:}
    \begin{itemize}
        \item \texttt{python evaluate.py} $\longrightarrow$ Reproduces Validation $\text{PQ} = 0.4618$.
        \item \texttt{python inference.py} $\longrightarrow$ Generates official champion \texttt{submission.csv} ($\text{LB} = 0.39$).
    \end{itemize}
\end{itemize}

\section{Conclusion}
In this work, we presented MAGFiLO 1.0, an end-to-end decoupled framework for full-disk solar filament instance segmentation. By combining analytical limb-darkening compensation, a high-resolution macro scout detector, dual native-resolution square crop U-Nets with Centerline Dice loss, and solar physics-inspired medial-axis collinear spine reconnection, our approach addresses scale disparity, limb attenuation, and topological flux rope fragmentation. Achieving a Validation PQ of 0.4618 and a Public Leaderboard score of 0.39, MAGFiLO 1.0 provides an accurate and operational methodology for automated solar feature cataloging and space weather monitoring.

\begin{thebibliography}{00}
\bibitem{parenti2014} S. Parenti, ``Solar Prominences: Observations,'' \textit{Living Reviews in Solar Physics}, vol. 11, no. 1, pp. 1--95, 2014.
\bibitem{kirillov2019panoptic} A. Kirillov, K. He, R. Girshick, C. Rother, and P. Doll{\'a}r, ``Panoptic Segmentation,'' in \textit{Proc. IEEE/CVF Conf. Comput. Vis. Pattern Recognit. (CVPR)}, 2019, pp. 9404--9413.
\bibitem{glenn2024yolo11} G. Jocher and J. Qiu, ``Ultralytics YOLO11,'' \textit{GitHub Repository}, 2024. [Online]. Available: \url{https://github.com/ultralytics/ultralytics}.
\bibitem{shit2021cldice} S. Shit, J. C. Paetzold, I. Sekuboyina, et al., ``clDice - a Novel Topology-Preserving Loss Function for Tubular Structure Segmentation,'' in \textit{Proc. IEEE/CVF Conf. Comput. Vis. Pattern Recognit. (CVPR)}, 2021, pp. 16560--16569.
\bibitem{bernasconi2005} P. N. Bernasconi, D. M. Rust, and D. Hakim, ``Advanced Automated Solar Filament Detection and Characterization Code: Description, Performance, and Results,'' \textit{Solar Physics}, vol. 228, no. 1, pp. 97--115, 2005.
\end{thebibliography}

\end{document}
